\begin{abstract}

Single-cell proteomics (SCP) generates protein abundance measurements across
hundreds to thousands of individual cells, offering unprecedented resolution to
study cellular heterogeneity. However, existing differential abundance (DA)
methods are limited to detecting shifts in mean expression, leaving biologically
relevant differences in shape undetected. Indeed, the specific power of SCP is
to identify differences between individual cells in a population, which are
typically only found as shape differences rather than in mean expression.
We therefore present \texttt{gamdid} (generalized additive
models for differential distributions), a novel statistical framework and R
package for differential distribution (DD) analysis in SCP data. \texttt{gamdid}
is based on generalized additive models (GAMs) to flexibly model heterogeneous
distributions, perform inference and provide interpretable visualizations.
Through semi-synthetic benchmarking on two SCP datasets, \texttt{gamdid}
demonstrates conservative false discovery rate (FDR) control
and substantially outperforms competing methods for differences in shape,
while achieving comparable performance for mean shifts. A spike-in case study
further demonstrates the utility of \texttt{gamdid} and its interpretable
visualization. Uniquely among DD methods, \texttt{gamdid} supports omnibus
testing across more than two groups, with post-hoc pairwise comparisons via
stagewise testing, and is specifically tailored for proteomics abundance data.

\end{abstract}
